% ============================================================
% Section 2: Background and Related Work
% Every work characterized here has been read. Sources:
%   ma2025understanding  - full text incl. Appendix D
%   xie2025realm         - full text
%   xue2023approxabft    - abstract + description in xie2025realm Sec II-C
%   huang1984abft, higham2002accuracy - standard references
% ============================================================

\section{Background and Related Work}
\label{sec:related}

\subsection{Silent data corruption}

Silent data corruption is computation that completes and is wrong. Google
traced silent corrupt execution errors to manufacturing defects in CPU cores
and reported that the detected incidence is higher than software engineers
expect \citep{hochschild2021cores}; Meta reported corruption at fleet scale
\citep{dixit2021silent}. Error-correcting codes protect memory and do not
address corruption arising inside the datapath, which is what makes detection
during computation necessary at all.

\subsection{Checksum ABFT}

\citet{huang1984abft} introduced algorithm-based fault tolerance for matrix
operations: augment a product with checksums and verify an invariant that holds
in exact arithmetic. In the classical two-sided form, $W$ and $X$ are augmented
with a checksum row $e^{\top}W$ and column $Xe$. Lighter one-sided variants
compute only $e^{\top}WX$ or $WXe$, and the matrix sum deviation
$e^{\top}WXe$ reduces the check to a single scalar
\citep[surveyed in][Sec.~II-C]{xie2025realm}. The detector studied here is the
one-sided form, compared per row rather than reduced to a scalar.

Whatever the form, the detector needs a tolerance, because finite-precision
rounding makes the two paths disagree even when the hardware is correct. How
that tolerance is chosen is where the literature divides, and it is the axis
this paper sits on.

\subsection{Thresholds from numerical bounds, and where they stop working}
\label{sec:related:bounds}

The first approach derives the tolerance from floating-point error analysis.
\citet{smith2015abft} give the formulation used as a detection rule in
subsequent work: report corruption when
$\lVert Cw - A(Bw) \rVert_\infty > \tau \lVert A \rVert_\infty
\lVert B \rVert_\infty$ with $\tau = Ku$ and $u$ the unit roundoff. The bound
behind $\tau$ requires $Ku \leq 0.01$ \citep[Eq.~2.7.11]{golub2013matrix}.

\citet{ma2025understanding} state the consequence explicitly. For an embedding
dimension of $4096$: at float32, $Ku \approx 0.0005$ and the condition holds; at
float16, $Ku \approx 4$; at bfloat16, $Ku \approx 32$. Past float32 the
condition fails, and the threshold is not conservative but undefined. They ran
their ABFT evaluation at float32 for this reason, report that running the same
check at bfloat16 produced false positives on nodes that had passed production
margin tests, and name extending ABFT to low-precision datatypes as open work.

That is the regime this paper occupies: the same contraction length of $4096$,
at bfloat16, where $Ku \approx 32$ and no principled threshold is available.

\subsection{Thresholds from model resilience}
\label{sec:related:resilience}

The second approach abandons the numerical bound and tunes the tolerance
against how much error the model tolerates before its task performance
degrades. \citet{xue2023approxabft} relax the strict comparison so that
recovery is invoked only when computing errors are large enough to matter,
reporting reduced computing overhead and improved accuracy against a strict
ABFT baseline. \citet{xie2025realm} extend this with a statistical criterion
for LLM inference, characterizing resilience across network components and
separating error magnitude from error frequency: they observe that components
followed by normalization are markedly more sensitive, because a single large
error becomes an outlier that skews the mean and variance the normalization
depends on. Their thresholds are fitted parameters, set empirically to allow a
perplexity increase of $0.3$ and an accuracy decrease of $0.5\%$.

This is calibration, but against a different quantity than ours, and the
distinction is not incidental. Both works operate on quantized accelerator
datapaths; \citet{xie2025realm} inject into INT32 GEMM results with INT8
inputs. In integer arithmetic the checksum identity is exact. There is no
rounding, no noise floor, and no false positives to trade against. Their
thresholds exist to avoid unnecessary \emph{recovery}, not to avoid spurious
\emph{detection}. In floating point the opposite constraint binds: the
threshold is bounded below by the rate at which clean hardware trips it, and
that rate has to come from somewhere.

\subsection{What is missing}

Neither family supplies a threshold for floating-point GEMM at a precision
where the classical bound has gone vacuous. The numerical-bound line stops at
float32 by its own stated condition. The model-resilience line calibrates
against task degradation in settings where the false-positive problem does not
arise. We are not aware of prior work that calibrates a checksum ABFT
threshold empirically at bfloat16 against a stated per-GEMM false-positive
rate, which is what this paper supplies, together with a measurement of what
the resulting detector fails to catch.

\subsection{Two negative results, from opposite directions}

\citet{ma2025understanding} also report that checksum ABFT failed to flag
silent corruption on fourteen of fifteen unhealthy production nodes, and
attribute this to real perturbations being small enough relative to the matrix
norm to fall inside floating-point error bounds. That is a negative result
about the detector, obtained from genuinely faulty hardware, at float32.

This paper reaches a compatible conclusion by a different route: controlled
injection at bfloat16, against an empirically calibrated threshold, with the
blind spot measured per bit class rather than inferred. The two results share
no method, no precision, and no corruption source. They agree that the case for
treating checksum ABFT as a complete SDC detector is weak.

\subsection{Fault injection and its abstraction level}

Fault injection can be performed at the algorithm level, the instruction level,
or the micro-architectural level, and the choice bounds what the results
support. Instruction-level studies of LLM inference on GPUs exist
\citep{chai2026llm}. Algorithm-level injection, used here and in the resilience
work above, does not reproduce instruction-level or micro-architectural
effects. We treat this as a scope boundary rather than a limitation to be
argued away; the model is specified in
Section~\ref{sec:method:perturbation}.
